\documentclass[10pt,twoside]{article}
\usepackage[utf8]{inputenc}
\usepackage{amsmath}
\usepackage{amsfonts}
\usepackage{amssymb}
\usepackage[spanish,es-noshorthands]{babel}
\usepackage[T1]{fontenc}
\usepackage{lmodern}
\usepackage{graphicx,hyperref}
\usepackage{multirow} 
\usepackage{tikz,pgf}
\usepackage{multicol}
\usepackage{subfig}
\usepackage[papersize={6.5in,8.5in},width=5.5in,height=7in]{geometry}
\usepackage{fancyhdr}
\pagestyle{fancy}
\fancyhead[LE]{\includegraphics[height=12pt]{Images/logo-colegio.png} Aritm\'etica $6^{\circ}$}
\fancyhead[RE]{}
\fancyhead[RO]{\textit{Germ\'an Avenda\~no Ram\'irez, Lic. U.D., M.Sc. U.N.}}
\fancyhead[LO]{}

\author{Germ\'an Avenda\~no Ram\'irez, Lic. U.D., M.Sc. U.N.}
\title{\begin{minipage}{.2\textwidth}
\includegraphics[height=1.75cm]{Images/logo-colegio.png}\end{minipage}
\begin{minipage}{.55\textwidth}
\begin{center}
Taller 07, Multiplicación en $\mathbb{N}$ \\
Aritm\'etica $6^{\circ}$
\end{center}
\end{minipage}\hfill
\begin{minipage}{.2\textwidth}
\includegraphics[height=1.75cm]{Images/logo-sed.png} 
\end{minipage}}
\date{}
\begin{document}
\maketitle
Nombre: \hrulefill Curso: \underline{\hspace*{44pt}} Fecha: \underline{\hspace*{2.5cm}}
\section*{Multiplicaci\'on en $\mathbb{N}$}
\subsection*{Acitividad 1}
Realice en forma individual en su cuaderno lo siguiente:
\begin{itemize}
 \item Tome una hoja de papel. D\'oblela de manera que queden, bien 4 filas y 8 columnas o bien, 8 filas y 4 columnas as\'i:
\begin{center}
\begin{tikzpicture}
\draw[help lines](0,0)grid(8,4);
\node[above] at (0.5,4){Columna 1};
\node[left]at(0,3.5){Fila 1};
\end{tikzpicture}
\end{center}
\item Responda las siguientes preguntas:
\begin{itemize}
 \item ?`En cu\'antas partes queda dividido el papel?
\item ?`Cu\'antos cuadrados tiene cada columna?
\item ?`Cu\'antos cuadrados tiene cada fila?
\item ?`Cu\'anto es 8 veces 4?, es decir, $4+4+4+4+4+4+4+4$
\item ?`Cu\'anto es 4 veces 8?, es decir, $8+8+8+8$.
\item ?`C\'omo se escribe abreviadamente 4 veces 8?, ?`8 veces 4?
\item ?`Qu\'e resultado se obtiene?
\end{itemize}
\item Recuerda:
\end{itemize}
La operación, que es una suma abreviada de sumandos iguales, se llama MULTIPLICACIÓN.
La multiplicación entre dos números naturales $a$ y $b$, se simboliza así:
\[a\cdot b \qquad \text{ó} \qquad a\times b, \qquad \mbox{8 veces 4}=8\cdot 4=8\times 4=4+4+4+4+4+4+4+4=32\]
El punto $\cdot$ y el signo $\times$ indican multiplicación. Cada término que interviene en la
operación se llama FACTOR. El número que se repite se llama MULTIPLICANDO y
el número de veces que el sumando se repite se llama MULTIPLICADOR
\begin{center}
\begin{tabular}{ccccc}
8 & $\times$ & 4 & = & 32 \\ 
Multiplicando &  & Multiplicador &  & Producto \\ 
\multicolumn{3}{c}{Factores}   &  &  \\ 
\end{tabular} 
\end{center}
\section*{Apliquemos las propiedades de la multiplicaci\'{o}n}
\subsection*{Actividad 2}
realizamos las siguientes operaciones y sacamos conclusiones.
\begin{itemize}
\item Respondamos en el cuaderno:
\begin{center}
\begin{tabular}{ll}
\multirow{2}{*}{\fbox{$2\times5=10$}}  & ¿Qué clase de números son el 2 y el 5? \\ 
 & ¿Qué clase de número es el 10? \\ 
\end{tabular} 
\end{center}
\begin{center}
\begin{tabular}{ll}
\multirow{2}{*}{\fbox{$3\times 4=12$}} & ¿Qué clase de números son el 3 y el 4? \\ 
 & ¿Qué clase de número es el 12? \\ 
\end{tabular} 
\end{center}
¿Qué clase de número es el producto de dos números naturales?
\item En el cuaderno realizamos las siguientes multiplicaciones:
\begin{itemize}
\begin{multicols}{2}
\item \begin{tabular}{|c|}
\hline 
$8\times 4=$? \\ 
\hline 
$4\times 8=$? \\ 
\hline 
\end{tabular} 
\item 
\begin{tabular}{|c|}
\hline 
$3\times 7=$? \\ 
\hline 
$7\times 3=$? \\ 
\hline 
\end{tabular}
\item
\begin{tabular}{|c|}
\hline 
$9\times 4=$? \\ 
\hline 
$4\times 9=$? \\ 
\hline 
\end{tabular}
\item 
\begin{tabular}{|c|}
\hline 
$5\times 1=$? \\ 
\hline 
$1\times 5=$? \\ 
\hline 
\end{tabular}  
\end{multicols}
¿Qué podemos concluir?
\end{itemize}
\item En el cuaderno realizamos las siguientes multiplicaciones:
\begin{itemize}
\begin{multicols}{2}
\item 
\begin{tabular}{|c|}
\hline 
$(4\times 2)\times 3=$? \\ 
\hline 
$4\times (2\times 3)=$? \\ 
\hline 
\end{tabular}
\item 
\begin{tabular}{|c|}
\hline 
$(3\times 2)\times 5=$? \\ 
\hline 
$3\times (2\times 5)=$? \\ 
\hline 
\end{tabular} 
\item 
\begin{tabular}{|c|}
\hline 
$(6\times 2)\times 3=$? \\ 
\hline 
$6\times (2\times 3)=$? \\ 
\hline 
\end{tabular} 
\item
\begin{tabular}{|c|}
\hline 
$(3\times 4)\times 3=$? \\ 
\hline 
$3\times (4\times 3)=$? \\ 
\hline 
\end{tabular} 
\end{multicols}
\end{itemize}
¿Qué conclusiones podemos sacar?
\item En el cuaderno. Contemos los puntos.

\begin{minipage}{.5\textwidth}
\begin{tikzpicture}
\draw[help lines] (0,0) grid (6,1);
\filldraw (.5,.5) circle(.5ex) (1.5,.5) circle (.5ex) (2.5,.5)circle(.5ex) (3.5,.5)circle(.5ex) (4.5,.5)circle(.5ex) (5.5,.5)circle(.5ex);
\end{tikzpicture}
\end{minipage}\hfill
\begin{minipage}{.4\textwidth}
6 veces 1 = 6 \\$1+1+1+1+1+1=6\times 1=6$
\end{minipage}

\begin{minipage}{.3\textwidth}
\begin{tikzpicture}
\draw (0,0) rectangle (1,3);
\filldraw (.5,.25)circle(.5ex) (.5,.75)circle(.5ex) (.5,1.25)circle(.5ex) (.5,1.75)circle(.5ex) (.5,2.25)circle(.5ex) (.5,2.75)circle(.5ex);
\end{tikzpicture}
\end{minipage}\hfill
\begin{minipage}{.6\textwidth}
Una vez seis\\
$1\times 6=$?
\end{minipage}
\end{itemize}
\end{document}
